\documentclass[11pt,a4paper]{article}

% ── Packages ──────────────────────────────────────────────────────────────────
\usepackage[top=2.8cm, bottom=2.8cm, left=3cm, right=3cm]{geometry}
\usepackage[T1]{fontenc}
\usepackage[utf8]{inputenc}
\usepackage{lmodern}
\usepackage{microtype}
\usepackage[colorlinks=true, linkcolor=linkblue, urlcolor=accent,
            pdftitle={L1-10 Distributed Systems -- System Design Foundations},
            bookmarks=true]{hyperref}
\usepackage{xcolor}
\usepackage{titlesec}
\usepackage{enumitem}
\usepackage{booktabs}
\usepackage{tabularx}
\usepackage{tcolorbox}
\usepackage{fancyhdr}
\usepackage{amsmath}
\usepackage{amssymb}
\usepackage{parskip}
\usepackage{mdframed}
\usepackage{graphicx}
\usepackage{tikz}
\usepackage{setspace}
\usepackage{soul}
\usepackage{array}
\usepackage{longtable}

% ── Colors ────────────────────────────────────────────────────────────────────
\definecolor{primary}{HTML}{1A1A2E}
\definecolor{accent}{HTML}{C0392B}
\definecolor{highlight}{HTML}{1A5276}
\definecolor{linkblue}{HTML}{154360}
\definecolor{lightbg}{HTML}{F4F6F7}
\definecolor{quizbg}{HTML}{EBF5FB}
\definecolor{termbg}{HTML}{EAF2FF}
\definecolor{curiobg}{HTML}{F5EEF8}
\definecolor{greenbg}{HTML}{EAFAF1}
\definecolor{notebg}{HTML}{FDF2E9}
\definecolor{accentlight}{HTML}{FADBD8}

% ── Section styling ───────────────────────────────────────────────────────────
\titleformat{\section}{\large\bfseries\color{highlight}}{\thesection.}{0.6em}{}[\vspace{2pt}\color{highlight!60}\rule{\linewidth}{0.6pt}]
\titleformat{\subsection}{\normalsize\bfseries\color{primary}}{\thesubsection}{0.5em}{}
\titleformat{\subsubsection}{\normalsize\itshape\color{primary!80}}{}{0em}{}
\titlespacing{\section}{0pt}{18pt}{8pt}
\titlespacing{\subsection}{0pt}{12pt}{4pt}

% ── Header/Footer ─────────────────────────────────────────────────────────────
\pagestyle{fancy}
\fancyhf{}
\fancyhead[L]{\small\color{highlight}\textbf{L1-10 --- Distributed Systems Fundamentals}}
\fancyhead[R]{\small\color{primary!70}System Design Interview Prep}
\fancyfoot[C]{\small\color{primary!60}\thepage}
\renewcommand{\headrulewidth}{0.4pt}
\renewcommand{\headrule}{\hbox to\headwidth{\color{highlight!40}\leaders\hrule height \headrulewidth\hfill}}

% ── Box environments ──────────────────────────────────────────────────────────
\tcbuselibrary{skins, breakable, hooks}

\newtcolorbox{termbox}[1]{enhanced,breakable,colback=termbg,colframe=highlight,coltitle=white,fonttitle=\bfseries\small,title={#1},left=8pt,right=8pt,top=6pt,bottom=6pt,attach boxed title to top left={yshift=-2.5mm,xshift=6mm},boxed title style={colback=highlight,sharp corners,left=4pt,right=4pt}}

\newtcolorbox{industrybox}[1]{enhanced,breakable,colback=greenbg,colframe=highlight!50!black,coltitle=white,fonttitle=\bfseries\small,title={Industry Data: #1},left=8pt,right=8pt,top=6pt,bottom=6pt,attach boxed title to top left={yshift=-2.5mm,xshift=6mm},boxed title style={colback=highlight!60!black,sharp corners,left=4pt,right=4pt}}

\newtcolorbox{trapbox}[1]{enhanced,colback=accentlight,colframe=accent,coltitle=white,fonttitle=\bfseries\small,title={Interview Trap: #1},left=8pt,right=8pt,top=6pt,bottom=6pt,attach boxed title to top left={yshift=-2.5mm,xshift=6mm},boxed title style={colback=accent,sharp corners,left=4pt,right=4pt}}

\newtcolorbox{thinkbox}{enhanced,colback=notebg,colframe=accent!60,left=8pt,right=8pt,top=6pt,bottom=6pt,leftrule=3pt,rightrule=0pt,toprule=0pt,bottomrule=0pt}

\newtcolorbox{curiobox}[1]{enhanced,breakable,colback=curiobg,colframe=primary!60,coltitle=white,fonttitle=\bfseries\small,title={Q: #1},left=8pt,right=8pt,top=6pt,bottom=6pt,attach boxed title to top left={yshift=-2.5mm,xshift=6mm},boxed title style={colback=primary!70,sharp corners,left=4pt,right=4pt}}

\newtcolorbox{quizbox}{enhanced,breakable,colback=quizbg,colframe=highlight,coltitle=white,fonttitle=\bfseries,title={Self-Quiz},left=10pt,right=10pt,top=8pt,bottom=8pt,attach boxed title to top left={yshift=-2.5mm,xshift=6mm},boxed title style={colback=highlight,sharp corners,left=4pt,right=4pt}}

\newtcolorbox{answerbox}{enhanced,breakable,colback=lightbg,colframe=primary!40,coltitle=white,fonttitle=\bfseries\small,title={Model Answers},left=10pt,right=10pt,top=8pt,bottom=8pt,attach boxed title to top left={yshift=-2.5mm,xshift=6mm},boxed title style={colback=primary!60,sharp corners,left=4pt,right=4pt}}

\newtcolorbox{insightbox}{enhanced,colback=primary!5,colframe=primary,left=8pt,right=8pt,top=6pt,bottom=6pt,leftrule=4pt,rightrule=0.4pt,toprule=0.4pt,bottomrule=0.4pt}

\setstretch{1.15}

% ─────────────────────────────────────────────────────────────────────────────
\begin{document}

% ── Title Page ────────────────────────────────────────────────────────────────
\begin{titlepage}
\newgeometry{top=3cm, bottom=3cm, left=3.5cm, right=3.5cm}
\thispagestyle{empty}

\vspace*{1.5cm}

{\color{highlight!60}\rule{\linewidth}{1.2pt}}
\vspace{0.4cm}

{\fontsize{9}{11}\selectfont\color{primary!60}\texttt{LESSON L1-10 \quad|\quad LAYER 1 --- FOUNDATIONS \quad|\quad SYSTEM DESIGN INTERVIEW PREP}}

\vspace{0.6cm}

{\fontsize{30}{36}\selectfont\bfseries\color{primary}%
Distributed Systems\\
\color{highlight}Fundamentals}

\vspace{0.4cm}

{\fontsize{14}{18}\selectfont\color{primary!75}\itshape%
CAP Theorem, Consensus, Consistency Models,\\
and the Tradeoffs That Define Every Real System}

\vspace{0.3cm}
{\color{highlight!60}\rule{\linewidth}{0.6pt}}

\vspace{1.0cm}

{\small\color{primary!65}%
\textbf{Edition:} 2025 \quad|\quad
\textbf{Data Sources:} Amazon Dynamo paper (2007), Google Spanner paper (2012/2017 TrueTime update),\\
CockroachDB engineering blog (2023), etcd project documentation (2024), Cloudflare engineering blog (2024),\\
Jepsen distributed systems analysis reports (2022--2024), Martin Kleppmann \textit{Designing Data-Intensive Applications} (2017)
}

\vspace{1.2cm}

\begin{insightbox}
\textbf{Why this lesson unlocks every tradeoff conversation in your interviews.}

Every system design interview reaches a moment where the interviewer asks: ``What happens if one of your database nodes goes down?'' or ``How do you handle a network partition between your primary and replica?'' or ``Would you use Cassandra or HBase here?'' Without a precise mental model of CAP theorem, consistency guarantees, and consensus algorithms, you will give vague answers like ``we use eventual consistency'' without being able to explain what that actually means, when it breaks, or why you chose it.

This lesson gives you the conceptual skeleton that sits beneath all distributed database choices. Whether you are designing a notification system, a payment ledger, a feed ranker, or a chat service --- the tradeoffs are the same, and CAP is always in the room. After this lesson, you will be able to name your consistency requirement before naming your database, which is the exact order of reasoning that separates senior engineers from junior ones in Apple ICT3-level interviews.

\textit{Note: In session\_001 (Notification System, 2026-04-15), the critical failure was jumping to database choices before articulating delivery semantics. Distributed systems fundamentals is the direct fix for that reasoning pattern.}
\end{insightbox}

\vfill
{\color{highlight!60}\rule{\linewidth}{0.6pt}}
{\small\color{primary!50} System Design Interview Prep Series \hfill L1-10 of 14}

\restoregeometry
\end{titlepage}

% ── Table of Contents ─────────────────────────────────────────────────────────
\newpage
\tableofcontents
\newpage

% ─────────────────────────────────────────────────────────────────────────────
\section{Why This Exists: The Single-Machine Illusion}
% ─────────────────────────────────────────────────────────────────────────────

For most of computing history, ``a server'' meant one machine. It had one CPU, one block of RAM, one disk. When you wrote to that disk, the write either succeeded or it failed. When you read data back, you got exactly what you wrote. If you wanted to know the time, you called the system clock and got a single authoritative answer. The state of your entire application lived in one place, and coordination was free --- a single mutex or a transaction was all you needed.

This model worked remarkably well until the mid-2000s, when internet applications began to scale beyond what any single machine could handle. Google's web index could not fit on one server. Amazon's product catalog could not serve millions of concurrent requests from a single machine. The solution was obvious in retrospect: spread the data and computation across many machines. But this introduced a problem that no one had fully reckoned with before.

When you have multiple machines, the network between them becomes an actor in your system. And networks are unreliable. Packets get dropped. Links go down. Routers fail. In a single-machine system, the CPU talks to RAM via a bus that is effectively perfect; failures are catastrophic (the machine crashes) rather than partial. But in a network, you get something far more insidious: \textit{partial failure}. Node A can talk to Node B, but not to Node C. Node C is alive and serving requests from clients, but Node B cannot reach it. Who has the authoritative version of the data? Who should accept writes?

This is the fundamental problem of distributed systems, and every concept in this lesson --- CAP theorem, consistency models, vector clocks, Raft, split-brain --- is a different angle on the same underlying challenge: how do you build a coherent system when the components can independently fail and the communication between them is unreliable?

\begin{thinkbox}
\textbf{Why not just use faster, more reliable networks?} Google, Amazon, and Meta operate their own fiber networks inside their data centers. Switch-to-switch latency is under a millisecond. Packet loss is measured in parts per million. And yet they still design all their distributed systems assuming the network will fail, because the mathematics of large-scale systems guarantees it: with 10,000 machines and a per-machine failure rate of 0.1\% per day, you statistically expect 10 failures every single day. Reliability engineering at scale means designing for failure as a steady state, not as an exceptional event.
\end{thinkbox}

% ─────────────────────────────────────────────────────────────────────────────
\section{CAP Theorem: What It Really Says (and What It Doesn't)}
% ─────────────────────────────────────────────────────────────────────────────

In 2000, Eric Brewer gave a keynote at PODC (Principles of Distributed Computing) proposing what became known as the CAP conjecture. In 2002, Gilbert and Lynch formally proved it as a theorem. The theorem states that a distributed data store can guarantee at most two of three properties simultaneously:

\begin{termbox}{CAP Theorem --- The Three Properties}
\textbf{Consistency (C):} Every read receives the most recent write or an error. Not ``eventually'' the most recent write --- the actual most recent write, immediately. This is sometimes called \textit{linearizability} or \textit{strong consistency}.

\textbf{Availability (A):} Every request (read or write) receives a response --- not an error, not a timeout. The system is always willing to respond, even if it cannot guarantee the response contains the latest data.

\textbf{Partition Tolerance (P):} The system continues to operate even when network partitions (dropped or delayed messages between nodes) occur.
\end{termbox}

Here is the single most important insight about CAP theorem, and the one that most candidates get wrong in interviews: \textbf{partition tolerance is not optional}. In any real distributed system deployed across more than one machine, network partitions \textit{will} happen. They are not edge cases; they are a normal operating condition. Therefore, the real choice in CAP theorem is not which two properties to pick --- it is which property you sacrifice \textit{during a partition}: consistency or availability.

When a partition occurs, the two sides of the partition can no longer communicate. If a write arrives on one side, the other side cannot know about it. You now have two choices. You can refuse to serve requests from either side until the partition heals (sacrificing availability, preserving consistency). Or you can serve requests from both sides, accepting that the two sides may diverge (preserving availability, sacrificing consistency). There is no third option. This is the fundamental tension.

\begin{thinkbox}
\textbf{What does ``consistency'' mean in CAP vs ACID?} These are completely different things and confusing them is one of the most common errors in distributed systems discussions. ACID consistency means the database always moves between valid states according to your schema constraints and integrity rules --- it is about application correctness. CAP consistency means every node in the distributed system returns the same value for the same key at the same logical time --- it is about synchronization across nodes. When you say ``I need a consistent database,'' an interviewer does not know which you mean. Be precise.
\end{thinkbox}

\subsection{Cassandra (AP) versus HBase (CP): A Real Example}

Apache Cassandra is the canonical example of an AP system. It was designed by Facebook in 2007 (released open source in 2008) to power the inbox search feature, which needed to handle billions of messages with high write throughput. The design choice was explicit: during a partition, Cassandra will continue accepting writes on all nodes. Different nodes may have different versions of the same row. When the partition heals, Cassandra uses a mechanism called ``last-write-wins'' (based on timestamps) or read repair to converge the data. The cost is that a client reading during a partition may get stale data.

Apache HBase, built on top of Hadoop HDFS, makes the opposite choice. HBase uses a single ``master'' architecture with a strong leader for each region of data. During a partition that isolates the region server from the master, HBase will stop accepting writes for that region rather than risk creating inconsistency. Clients see errors rather than stale data. This makes HBase a CP system: consistent and partition-tolerant, but not available during partitions.

\begin{industrybox}{Facebook / Apache Cassandra --- Availability at Scale, 2008}
Facebook's inbox search had to handle roughly 50 billion messages at launch. The team explicitly chose availability over consistency: if a user's inbox shows a message as read when it is actually unread (because the partition has not healed yet), that is a tolerable inconsistency. If the inbox fails to load entirely because the system is waiting for consistency, that is an intolerable user-facing failure. The CAP choice was driven directly by the business requirement. This is the correct reasoning pattern: start from the user-facing consequence, then pick your tradeoff.
\end{industrybox}

\subsection{The Common Misconceptions}

The most frequent interview mistake is treating CAP as a permanent architectural label rather than a dynamic property. A system is not ``a CP system'' or ``an AP system'' in all situations --- it is a system that chooses C over A (or vice versa) \textit{during a partition}. When there is no partition (which is the vast majority of time in a healthy network), the system can deliver both consistency and availability simultaneously. CAP only forces a choice under partition.

\begin{trapbox}{Treating CAP as a Binary, Static Label}
Saying ``Cassandra is an AP database'' is a shorthand, not a complete description. More precisely: Cassandra, by default, sacrifices consistency in favor of availability during network partitions. But Cassandra also allows you to tune its consistency level per-read and per-write using its quorum settings (ONE, QUORUM, ALL). If you configure Cassandra with \texttt{QUORUM} reads and \texttt{QUORUM} writes on a cluster of $N$ nodes, you get stronger consistency guarantees --- at the cost of higher latency and reduced availability. The system is not fixed; the tradeoff is tunable. When an interviewer asks ``which database would you use?'', they are partly testing whether you know to ask \textit{what consistency level do we need?} before answering.
\end{trapbox}

% ─────────────────────────────────────────────────────────────────────────────
\section{PACELC: The Better Model for Real Systems}
% ─────────────────────────────────────────────────────────────────────────────

CAP theorem describes behavior only during network partitions. But real systems are partitioned for only a small fraction of their operating lifetime. For the other 99.9\% of the time, when everything is working normally, CAP says nothing about performance. This is where PACELC fills the gap.

\begin{termbox}{PACELC --- The Extended Model (Daniel Abadi, 2012)}
PACELC stands for: \textbf{P}artition $\rightarrow$ \textbf{A}vailability vs. \textbf{C}onsistency; \textbf{E}lse $\rightarrow$ \textbf{L}atency vs. \textbf{C}onsistency.

The insight is that even without a partition, there is a fundamental tradeoff in distributed systems: to achieve strong consistency, every write must be acknowledged by multiple nodes before it is considered committed. Waiting for multiple nodes takes time. Therefore, strong consistency always costs latency. Alternatively, you can acknowledge a write to the client as soon as one node (the primary) has written it, and replicate asynchronously in the background. This reduces latency but means a read from a replica might get a stale value.

Every distributed database makes both choices: how to behave during partitions (the CA tradeoff from CAP), and how to balance latency against consistency in the normal case (the LC tradeoff in PACELC).
\end{termbox}

Google Spanner, released to the public as a paper in 2012 and available as Cloud Spanner since 2017, is classified as PA/EL in PACELC: during partitions it sacrifices availability (refuses requests rather than returning inconsistent data), and in the normal case it also prefers consistency over latency. Spanner achieves this using a globally synchronized clock called TrueTime, which uses GPS receivers and atomic clocks in each Google data center to bound clock uncertainty to within a few milliseconds. By knowing the worst-case clock skew, Spanner can assign globally consistent timestamps without needing distributed coordination for every transaction. In 2024, Google reported that Cloud Spanner serves customers including PayPal, Snap, and Deutsche Bank for workloads that require global ACID transactions at scale.

DynamoDB (Amazon's managed NoSQL service) is classified PA/EL in the default configuration: it favors availability over consistency during partitions, and it favors latency over consistency in the normal case (writes are acknowledged when the primary node commits them, before replicas confirm). But DynamoDB also supports ``strongly consistent reads'' as an option, at the cost of higher latency and cost --- shifting to PC/EC for those specific reads. This illustrates the key insight: PACELC is not a fixed label but a description of the default operating point, which can often be tuned.

\begin{industrybox}{Amazon DynamoDB --- Tunable Consistency, 2022}
Amazon's 2022 DynamoDB paper (published at USENIX ATC) describes how DynamoDB serves tens of millions of requests per second across millions of customers with 99.999\% availability. The paper notes that the vast majority of DynamoDB customers use eventually consistent reads (the default) and never need strong consistency, because their workloads tolerate brief inconsistency windows of under a second. The minority who need strong consistency pay a roughly 2x cost in latency and monetary cost. The lesson: choose your consistency model based on what the application actually requires, not on what sounds safest.
\end{industrybox}

% ─────────────────────────────────────────────────────────────────────────────
\section{The Consistency Models Spectrum}
% ─────────────────────────────────────────────────────────────────────────────

``Consistency'' is not a binary. There is a spectrum of guarantees, ranging from the strongest (linearizability) to the weakest (eventual consistency), and each point on the spectrum enables different performance characteristics.

\subsection{Linearizability (Strong Consistency)}

Linearizability is the gold standard of consistency. A system is linearizable if every operation appears to take effect instantaneously at some point between when the client sent the request and when it received the response, and all operations appear in a global total order consistent with real time. In plain English: every client sees the same state, and reads always reflect the latest write.

Achieving linearizability requires coordination. Every write must be confirmed by a quorum of nodes before the client gets a success response. Every read must either contact a quorum or read from a node that is guaranteed to have the latest data. This coordination costs latency --- typically one full round-trip of network communication for every operation. At the scale of a single data center with sub-millisecond latency, this cost is often acceptable. Across data centers with 50--200ms cross-region latency, it becomes a major bottleneck.

\begin{thinkbox}
\textbf{Is linearizability the same as serializability?} No, and this distinction matters in interviews. Linearizability is a property of single operations (reads and writes) across multiple nodes --- it is about the visibility of individual operations in real time. Serializability is a property of \textit{transactions} (multiple operations bundled together) --- it guarantees that concurrent transactions appear to execute in some serial order. A system can be serializable without being linearizable. A system that provides both is said to be \textit{strictly serializable} (also called \textit{external consistency}). Google Spanner and CockroachDB claim strict serializability, which is the strongest possible guarantee.
\end{thinkbox}

\subsection{Sequential Consistency}

Sequential consistency relaxes linearizability in one specific way: operations still appear in the same global order to all clients, but the order does not have to match real time. If Client A writes X=1 at time 10:00:00.000 and Client B writes X=2 at time 10:00:00.001, a sequentially consistent system is allowed to apply B's write before A's write --- but all clients will see the same order, whatever that order is.

This distinction sounds subtle but matters in practice. Sequential consistency is often achievable with lower latency than linearizability because you do not need a synchronized global clock or cross-node coordination for reads --- you only need to ensure that operations from a single client are applied in the order they were issued. Many database replication systems offer sequential consistency within a single client session.

\subsection{Causal Consistency}

Causal consistency captures only the relationships that actually matter for correctness: if event B was causally caused by event A (because B read a value that A wrote, or because a human did A and then did B), then all nodes must apply them in that order. Events that are causally independent can be applied in any order.

This model is more powerful than it might initially appear. Most real application logic is causal: a user posts a message (event A), then replies to it (event B). You need B to appear after A. But two different users posting unrelated messages have no causal relationship, so their relative order does not matter. Causal consistency captures this naturally, with much lower coordination overhead than linearizability or sequential consistency.

Amazon's Dynamo paper (2007) uses a mechanism called vector clocks to track causality, which we explore in Section~\ref{sec:vector-clocks}. The CRDTs (Conflict-free Replicated Data Types) used in systems like Riak are another approach to encoding causal relationships directly in the data structure.

\subsection{Eventual Consistency}

Eventual consistency is the weakest useful guarantee: if no new updates are made to a piece of data, all replicas will eventually converge to the same value. It says nothing about how long convergence takes or what a client will see during the convergence window.

This sounds weak, but it is precisely the right model for a large class of real-world applications. DNS is eventually consistent: when you update a DNS record, it takes minutes to hours to propagate to all resolvers worldwide, and during that window, different clients will resolve the same hostname to different IP addresses. This is completely acceptable because the application (loading a website) tolerates brief inconsistency far better than it would tolerate DNS lookups taking seconds to synchronize globally.

Shopping carts on Amazon's retail site are eventually consistent, as documented in the original Dynamo paper. If you add an item to your cart on your laptop and simultaneously add a different item on your phone, the two adds may temporarily be invisible to each other. The system will resolve this conflict (using the union of both adds, so you keep both items) and converge. The alternative --- a linearizable shopping cart that requires global coordination on every add --- would make checkout pages measurably slower for hundreds of millions of users, which is a much worse outcome.

\begin{industrybox}{Amazon Dynamo --- Eventual Consistency at Retail Scale, 2007}
The Amazon Dynamo paper, published at SOSP 2007, remains one of the most influential papers in distributed systems. It describes how Amazon's shopping cart, session state, and bestseller list were all moved to an eventually consistent store to achieve the availability and latency they needed at scale. The paper introduced the concept of ``always writeable'' --- the system never rejects a write, even during a partition. The key design insight: availability failures (a user cannot add an item to their cart) cost Amazon money immediately and visibly. Consistency failures (a cart shows a stale item count briefly) cost almost nothing in practice. Match the consistency model to the cost of the failure it prevents.
\end{industrybox}

% ─────────────────────────────────────────────────────────────────────────────
\section{Vector Clocks: Tracking Causality Across Nodes}
\label{sec:vector-clocks}
% ─────────────────────────────────────────────────────────────────────────────

When you have multiple nodes that can accept writes independently --- as in an AP system --- you need a way to determine which version of a value is newer, and whether two conflicting writes happened independently (concurrent) or one happened after the other (causally ordered). Wall-clock timestamps cannot do this reliably: clocks on different machines drift, and clock drift means two events with the same timestamp might have happened in either order, or one might have caused the other.

\begin{termbox}{Vector Clock --- Logical Causality Tracking}
A vector clock is a data structure that tracks causality in a distributed system without relying on synchronized physical clocks. Each node maintains a vector (array) of counters, one per node in the system. When node $i$ performs a local event, it increments its own counter in the vector. When node $i$ sends a message to node $j$, it includes its current vector. When node $j$ receives the message, it takes the element-wise maximum of its own vector and the received vector, then increments its own counter. Two events can be compared: event A \textit{causally precedes} event B if every counter in A's vector is $\leq$ the corresponding counter in B's vector (and at least one is strictly less). If neither event causally precedes the other, they are \textit{concurrent} --- they happened independently without knowledge of each other.
\end{termbox}

Consider three nodes: Node 1, Node 2, and Node 3. Initially all vectors are $[0, 0, 0]$. Node 1 writes a value, incrementing its vector to $[1, 0, 0]$ and attaching this to the value. Node 1 sends the value to Node 2. Node 2 updates its vector to $[\max(1,0), \max(0,0)+1, \max(0,0)] = [1, 1, 0]$. Meanwhile, Node 3, which did not receive Node 1's write, writes its own value with vector $[0, 0, 1]$. When we compare Node 2's version $[1, 1, 0]$ and Node 3's version $[0, 0, 1]$, neither dominates the other: Node 2's first counter is larger, but Node 3's third counter is larger. These writes are concurrent. A conflict resolution strategy must decide what to do --- for a shopping cart, the Dynamo approach is to take the union of both versions and present both to the client to reconcile.

The practical insight from vector clocks is subtle but important: they allow a distributed system to detect when two writes are genuinely in conflict (concurrent) versus when one write clearly supersedes another (causally ordered), without needing any coordination at write time. This is what makes Dynamo-style systems work: writes are always accepted, and conflicts are detected and resolved later using the causal history embedded in the vector clocks.

% ─────────────────────────────────────────────────────────────────────────────
\section{Raft: How Distributed Consensus Actually Works}
% ─────────────────────────────────────────────────────────────────────────────

For systems that require strong consistency --- that need a single authoritative version of truth across multiple nodes --- you need distributed consensus. Consensus means a cluster of nodes agreeing on a single value or sequence of operations, even in the presence of node failures. The classic algorithm is Paxos (Lamport, 1989), but Paxos is notoriously difficult to understand and implement correctly. In 2014, Diego Ongaro and John Ousterhout at Stanford published Raft, explicitly designed to be understandable. Raft is now the consensus algorithm of choice for most modern distributed systems.

\subsection{Leader Election}

Raft organizes time into \textit{terms}, each with a monotonically increasing term number. At any given time, each node is in one of three states: \textit{leader}, \textit{follower}, or \textit{candidate}. In steady state, there is exactly one leader per term, and all reads and writes go through the leader.

Nodes start as followers and expect to hear heartbeats from the leader. If a follower does not hear from the leader within an \textit{election timeout} (typically 150--300ms, randomized to prevent split votes), it becomes a candidate. It increments the term number, votes for itself, and sends RequestVote RPC messages to all other nodes. A node that has not yet voted in this term will vote for the candidate, provided the candidate's log is at least as up-to-date as the voter's log (the log freshness check prevents electing a node that has missed recent writes). If a candidate receives votes from a majority of nodes, it becomes the new leader and begins sending heartbeats to suppress further elections.

The randomized election timeout is clever. If all followers had the same timeout, they would all become candidates simultaneously after the leader failed, and no candidate would get a majority. By randomizing timeouts (each node picks a different delay from a range), one node will typically timeout first, become a candidate, and win the election before the others even try.

\begin{industrybox}{etcd in Kubernetes --- Raft for Cluster State, 2024}
etcd is a distributed key-value store that uses Raft for consensus. It is the brain of every Kubernetes cluster: it stores all cluster state, including which pods are running, which nodes are healthy, and what the desired configuration is. As of 2024, etcd is used by essentially every major Kubernetes deployment --- Google GKE, Amazon EKS, Azure AKS --- making Raft the consensus algorithm underpinning the infrastructure of most cloud-native applications. etcd documentation recommends a cluster of 3 or 5 nodes, because a 3-node cluster can tolerate 1 failure and a 5-node cluster can tolerate 2 failures (Raft requires a majority, meaning $(N+1)/2$ nodes must be healthy for the cluster to function). A single-node etcd is a common misunderstanding; without a majority requirement there is no consensus and no partition tolerance.
\end{industrybox}

\subsection{Log Replication}

Once a leader is elected, all writes go through it. The leader appends the write to its local log as a ``log entry'' and simultaneously sends AppendEntries RPC messages to all followers, asking them to append the same entry. A log entry is \textit{committed} once the leader has received acknowledgments from a majority of nodes (including itself). Committed means the entry will survive any future failures. The leader then applies the committed entry to its state machine and responds to the client.

This two-phase structure (append, then commit after majority acks) is what gives Raft its durability guarantee. A write acknowledged to the client is guaranteed to be present on a majority of nodes. Even if the leader then immediately crashes, a new election will produce a new leader that has this write in its log (because any node that can win an election must have a log at least as fresh as a majority, and a majority of nodes have the write). The client's write will not be lost.

\subsection{Safety Guarantee}

Raft's core safety property is: at most one leader per term. This is guaranteed by the election mechanism (a node can only be elected if a majority votes for it, and each node votes at most once per term). Combined with the log freshness requirement for voting, this ensures that any newly elected leader has all previously committed log entries. There can never be two nodes simultaneously believing they are the committed leader for the same term, which eliminates the possibility of conflicting committed writes.

\begin{industrybox}{CockroachDB and TiKV --- Raft for Distributed SQL, 2023}
CockroachDB (open-sourced by Cockroach Labs) and TiKV (the storage layer of TiDB, open-sourced by PingCAP) both use Raft to build distributed SQL databases. CockroachDB partitions data into ranges of approximately 64MB each, and each range is replicated via its own independent Raft group across 3 or 5 nodes. This means a CockroachDB cluster with data spread across 3 data centers can survive the complete loss of one data center without losing any committed data and with automatic failover, typically within seconds. The 2023 CockroachDB engineering blog reports that their largest customers run clusters with thousands of Raft groups, handling millions of transactions per day with cross-region ACID guarantees. This was effectively impossible before Raft made consensus practical to implement correctly at scale.
\end{industrybox}

% ─────────────────────────────────────────────────────────────────────────────
\section{Split-Brain: The Deadliest Distributed Systems Failure}
% ─────────────────────────────────────────────────────────────────────────────

Split-brain is what happens when a network partition creates two groups of nodes, each believing the other group has failed, and each proceeding to act as if it is the surviving primary. In a database context, this means two nodes simultaneously accepting writes to the same data, creating two divergent histories that must eventually be reconciled --- if they can be reconciled at all.

For data with natural conflict resolution (shopping carts where you take the union, counters where you take the sum), split-brain is recoverable. For data where writes conflict (two machines both accepting a payment from an account with insufficient funds for both, or two cluster controllers both scheduling a workload on the same resources), split-brain can cause data corruption or logical inconsistency that is very difficult to undo.

The primary defense against split-brain is the \textit{quorum} requirement: any node that cannot contact a majority of the cluster refuses to accept writes. This is exactly what Raft enforces: a Raft leader that loses contact with a majority of followers will step down and stop accepting writes. A new leader can only be elected if a majority of nodes are reachable and agree. Since a majority in a cluster of $N$ nodes requires more than $N/2$ nodes, two different majorities cannot simultaneously exist. Therefore, at most one subset of the cluster can elect a leader and accept writes at any given time.

\begin{thinkbox}
\textbf{What is STONITH and why do you need it if quorums prevent split-brain?} STONITH (Shoot The Other Node In The Head) is a fencing mechanism used in high-availability clusters, particularly for non-Raft systems like traditional Primary/Replica database clusters. In a Raft cluster, quorum mathematically prevents split-brain. But in older systems where a primary might not know it should step down (because it is still processing requests when it loses network connectivity), STONITH provides a hardware-level mechanism to power off the rogue primary node via a remote-controlled power switch or management interface. The name is deliberately dramatic: you literally kill the other node to ensure only one primary exists. Modern cloud environments use equivalent mechanisms (forcibly terminating an EC2 instance via the AWS control plane, for example) when a node needs to be fenced.
\end{thinkbox}

\begin{trapbox}{Assuming Replicas Prevent Split-Brain}
A common interview mistake is confusing replication with split-brain protection. Traditional primary-replica replication (the default setup in MySQL, PostgreSQL, Redis) does \textit{not} prevent split-brain. If the primary loses network connectivity and the monitoring system promotes a replica to primary, the original primary may still be alive and accepting writes from clients that can still reach it. Now you have two primaries. Quorum-based consensus (Raft, Paxos) is the correct solution. Alternatively, you can accept split-brain and build conflict resolution into your data model --- but that requires an explicit architectural decision, not an accident.
\end{trapbox}

% ─────────────────────────────────────────────────────────────────────────────
\section{Two-Phase Commit and Why It Fails at Scale}
% ─────────────────────────────────────────────────────────────────────────────

Before Raft became dominant, the standard mechanism for coordinating writes across multiple nodes in a distributed transaction was Two-Phase Commit (2PC). Understanding why 2PC is problematic is as important as understanding why Raft works, because you will encounter 2PC in legacy systems and in discussions of distributed transactions.

\begin{termbox}{Two-Phase Commit (2PC) --- The Protocol}
\textbf{Phase 1 --- Prepare:} A coordinator node sends a PREPARE message to all participant nodes. Each participant checks whether it can commit the transaction (validates constraints, acquires locks, writes to its WAL) and responds YES or NO. Participants that respond YES are in a ``prepared'' state --- they have promised to commit if instructed.

\textbf{Phase 2 --- Commit or Abort:} If all participants responded YES, the coordinator sends COMMIT to all participants. Each participant commits the transaction, releases locks, and acknowledges. If any participant responded NO, the coordinator sends ABORT to all participants, and each rolls back.
\end{termbox}

The protocol sounds reasonable until you consider what happens when the coordinator fails after Phase 1 completes but before Phase 2 begins. All participants are in the prepared state: they have locked the relevant rows and are waiting for the commit or abort instruction. But the coordinator is dead. The participants cannot proceed without the coordinator's decision, because committing without knowing whether all other participants also said YES would violate atomicity. But rolling back is also dangerous: if some participants (those that the coordinator reached before crashing) have already committed, rolling back the others would violate atomicity in the other direction.

The result is called the \textit{blocking problem}: participants are stuck holding locks, unable to proceed, until the coordinator recovers. In a distributed system with thousands of nodes, the probability that a coordinator fails in this narrow window is low but non-zero, and at scale, ``non-zero'' happens regularly. Systems using 2PC across high-volume workloads frequently experience transactions blocked for minutes during coordinator recovery, holding locks that block other transactions.

\begin{thinkbox}
\textbf{Does Google Spanner use 2PC? If so, why is it reliable?} Yes, Spanner uses 2PC for distributed transactions --- but with a crucial addition. Spanner's 2PC uses Paxos groups (equivalent to Raft) for each participant, not single nodes. The Paxos group for each participant can elect a new leader if the coordinator fails, allowing it to complete the in-flight transaction. The coordinator is also a Paxos group, not a single node. So the ``coordinator failure'' scenario that breaks classical 2PC is handled by Paxos recovery within the coordinator group. This is the pattern for making 2PC production-grade: replace every single-point-of-failure node with a consensus group. The cost is significantly higher complexity.
\end{thinkbox}

Modern distributed databases have largely replaced 2PC with consensus-based approaches. CockroachDB uses Raft without 2PC for most operations. TiDB uses Percolator (Google's 2PC variant that uses MVCC to avoid the blocking problem) for cross-shard transactions. The lesson: 2PC is conceptually important and still appears in legacy enterprise middleware, but it is not the right primitive for building new distributed systems.

% ─────────────────────────────────────────────────────────────────────────────
\section{Decision Framework: Choosing CP vs AP for Your System}
% ─────────────────────────────────────────────────────────────────────────────

The most important skill from this lesson is the ability to derive your consistency requirement from the business requirement, then select technology that matches. The following framework walks through that reasoning.

\textbf{Step 1: What is the cost of a stale read?} If a user reads data that is slightly out of date, what happens? For a social media like count (off by a few during a partition), nothing bad happens. For a financial balance check before a withdrawal (off by the amount of a concurrent deposit), fraud becomes possible. The higher the cost of a stale read, the stronger your consistency requirement.

\textbf{Step 2: What is the cost of a failed write?} If you refuse to accept a write during a partition, what is the business impact? For a system that records sensor telemetry, a brief write failure during a partition means you lose some data points but can resume. For a user trying to place an order, a write failure means a lost sale. The higher the cost of a write failure, the stronger your availability requirement.

\textbf{Step 3: Can conflicts be resolved automatically?} For counters, sets, and last-write-wins data, conflicts resulting from an AP design can be resolved mechanically. For sequential IDs, financial ledgers, and inventory counts where overselling is catastrophic, conflicts cannot be resolved after the fact --- you need CP design to prevent them.

\vspace{0.5em}
\noindent\textbf{Comparison Table:}

\begin{center}
\begin{tabular}{>{\bfseries}p{3.2cm} p{3.5cm} p{3.5cm} p{3.8cm}}
\toprule
\textbf{System Type} & \textbf{Consistency Need} & \textbf{Availability Need} & \textbf{Recommended Approach} \\
\midrule
Financial ledger / payments & Linearizable (no stale reads, no conflicts) & Medium (errors ok if data at risk) & CP: Raft-based DB (CockroachDB, Spanner) \\
\addlinespace
Inventory / stock count & Strong (oversell = loss) & High (lost sales are costly) & CP with short partition window; or optimistic locking \\
\addlinespace
User session / auth token & Causal (session coherent per user) & Very High (auth failure = logout) & Causal: consistent hashing + session affinity \\
\addlinespace
Social feed / likes / follower counts & Eventual (brief lag ok) & Very High (user sees nothing = bad) & AP: Cassandra, DynamoDB (eventual) \\
\addlinespace
Distributed configuration (k8s, service discovery) & Linearizable (wrong config = outage) & Medium (brief unavailability tolerable) & CP: etcd (Raft) \\
\addlinespace
Shopping cart & Eventual + conflict-free (union of adds) & Very High (cart failure = lost sale) & AP: Dynamo-style, CRDT-based resolution \\
\addlinespace
Distributed lock / leader election & Linearizable (exactly one holder) & Low (fine to fail if unsure) & CP: ZooKeeper, etcd \\
\bottomrule
\end{tabular}
\end{center}

% ─────────────────────────────────────────────────────────────────────────────
\section{Critical Numbers Reference}
% ─────────────────────────────────────────────────────────────────────────────

\begin{center}
\begin{tabular}{p{5.5cm} p{4.2cm} p{4.3cm}}
\toprule
\textbf{Metric} & \textbf{Value} & \textbf{Source / Year} \\
\midrule
Raft election timeout (typical) & 150--300 ms (randomized) & Raft paper, 2014 \\
\addlinespace
Raft heartbeat interval & 50--150 ms & Raft paper, 2014 \\
\addlinespace
etcd recommended cluster size & 3 or 5 nodes & etcd docs, 2024 \\
\addlinespace
etcd max recommended cluster size & 7 nodes & etcd docs, 2024 \\
\addlinespace
Raft min nodes for $k$ failures & $2k + 1$ nodes & Raft paper, 2014 \\
\addlinespace
CockroachDB range size & $\sim$64 MB per Raft range & CockroachDB docs, 2023 \\
\addlinespace
Google TrueTime clock uncertainty & $<$ 7 ms (typical $<$ 4 ms) & Spanner paper, 2012 \\
\addlinespace
Dynamo vector clock max size & Bounded by node count; Dynamo used $\sim$10 entries & Dynamo paper, 2007 \\
\addlinespace
Cross-DC network latency (US East to West) & 60--80 ms RTT & Cloudflare, 2023 \\
\addlinespace
Cross-DC network latency (US to Europe) & 100--150 ms RTT & Cloudflare, 2023 \\
\addlinespace
2PC blocking window (coordinator failure) & Duration of coordinator restart; seconds to minutes & Bernstein et al., 1987 \\
\addlinespace
Cassandra default replication factor & 3 & Cassandra docs, 2024 \\
\addlinespace
DynamoDB 99.999\% availability SLA & $<$ 5.3 min downtime per year & AWS SLA, 2024 \\
\bottomrule
\end{tabular}
\end{center}

% ─────────────────────────────────────────────────────────────────────────────
\section{System Design Application}
% ─────────────────────────────────────────────────────────────────────────────

\subsection{Application 1: Designing a Distributed Lock Service}

A distributed lock ensures that at most one client holds the lock at any given time, even across multiple nodes and in the presence of network partitions. This appears in system design interviews when you need leader election, rate limiting, or exclusive resource access.

The correct approach is a CP system. Availability does not help here: if you cannot contact a quorum, you must not grant the lock, because granting it without certainty risks two clients holding the lock simultaneously. etcd provides distributed locks via its \texttt{/v3/lock} API. The lock is implemented as a key with a lease (TTL). The client that creates the key wins the lock. If the client crashes without releasing the lock, the lease expires and the lock is released automatically. etcd guarantees that key creation is linearizable --- only one client can create a given key, and all other clients that simultaneously attempt it will fail. The lease is tracked via Raft, so if the etcd leader fails, the new leader inherits all active leases and their remaining TTLs.

The common interview mistake here is proposing Redis for distributed locking. Redis's SETNX command provides a primitive lock, but standard Redis is a single node. Redis Cluster does not provide a linearizable lock across nodes --- the Redlock algorithm attempts to address this using multiple independent Redis nodes, but Martin Kleppmann's 2016 analysis demonstrated that Redlock is not safe under all timing scenarios (clock drift and process pauses can cause two clients to simultaneously hold the lock). For production-grade distributed locking, use etcd or ZooKeeper.

\subsection{Application 2: Cross-Region Leader Election for a Stateful Service}

Suppose you have a stateful processing service (a payment processor, an order coordinator, a scheduling engine) that must have exactly one active leader at a time, even as you deploy it across three geographic regions for disaster recovery. If the leader region fails, a new leader must be elected within seconds.

The design: deploy one instance in each region. Use etcd (or ZooKeeper) as the consensus layer, also deployed across all three regions as a 3-node cluster (one node per region). Each service instance tries to acquire a lease in etcd on startup. The instance that wins the lease is the leader. All other instances watch the lease and stand by. If the leader crashes or the lease expires, a new election occurs within one Raft election timeout (150--300ms) plus the lease expiry time. The lease TTL should be set to balance rapid failover (short TTL) against false failovers due to temporary network hiccups (long TTL). Kubernetes uses this pattern for its own controller manager and scheduler, using a 15-second lease with a 5-second renewal interval.

\subsection{Application 3: Cross-Region Database Replication with a Consistency Guarantee}

For a globally deployed application (say, a note-taking app with users in the US, Europe, and Asia), you want reads and writes to be served from the nearest region, but you need all regions to have the same view of the data. This is a cross-region replication design.

If your application can tolerate eventual consistency (notes created on one device appear on another device within a few seconds), the pattern is: use a multi-region AP database (DynamoDB Global Tables, Cassandra multi-DC, or CockroachDB with high latency tolerance). Writes go to the nearest region and are asynchronously replicated to other regions. Reads from a given region see a locally-cached version that is eventually consistent with other regions. Conflict resolution for concurrent writes to the same document from different regions uses last-write-wins (timestamp-based) or CRDTs.

If your application requires strong consistency (a financial ledger where a user must see their balance immediately after a transfer, even from a different device in a different country), the pattern is fundamentally harder. Google Spanner achieves this by assigning timestamps using TrueTime and committing transactions with a delay equal to the maximum clock uncertainty ($\sim$7ms), ensuring all replicas have applied all earlier transactions before the current one is visible. The cost is that every write takes at least 7ms longer than it would without the global ordering guarantee. For sub-7ms write latency requirements and cross-region strong consistency, there is currently no solution that avoids this fundamental physics constraint.

% ─────────────────────────────────────────────────────────────────────────────
\section{Curiosity Corner}
% ─────────────────────────────────────────────────────────────────────────────

\begin{curiobox}{If partition tolerance is always required, why do some databases label themselves ``CA''?}
You will occasionally see systems (particularly older relational databases like traditional single-node PostgreSQL or MySQL) described as ``CA'' systems in CAP taxonomy. This is technically valid but practically misleading. A single-node database is not a distributed system --- it runs on one machine, so there is no network between components and therefore no possibility of a network partition in the CAP sense. It provides both consistency and availability trivially because the ``partition'' scenario simply does not apply. When you replicate that database to a second node (to achieve high availability or read scaling), you immediately re-enter the distributed system problem space and must choose between C and A during a partition. ``CA'' is therefore a label for non-distributed systems, not a third option in the CAP tradeoff for distributed databases.
\end{curiobox}

\begin{curiobox}{Can you get linearizability without Raft or Paxos?}
Technically yes, through other mechanisms. Timestamp ordering with synchronized clocks (as in Google Spanner's TrueTime) provides linearizability without a traditional consensus round per operation, because the globally synchronized clock provides a total order. Multi-Version Concurrency Control (MVCC) with sufficiently precise timestamps can achieve serializable isolation. However, clock synchronization at the sub-millisecond level requires specialized hardware (GPS, atomic clocks) and careful measurement of uncertainty bounds --- Google's solution is not easily replicable outside hyperscaler infrastructure. For most practical distributed systems without access to TrueTime-equivalent infrastructure, Raft or Paxos is the right approach. There are also novel approaches like Calvin (deterministic database), which uses a fixed global ordering layer to pre-assign transaction sequences without consensus rounds, but these are not yet mainstream.
\end{curiobox}

\begin{curiobox}{What is the difference between Raft and Multi-Paxos, and why does it matter?}
Raft and Multi-Paxos are functionally equivalent in what they guarantee --- both are consensus algorithms that provide linearizable log replication across a majority quorum. The differences are in design philosophy and implementation clarity. Multi-Paxos (the version of Paxos used in practice, not the original single-decree Paxos) was specified informally by Lamport, leaving many implementation details unspecified: how to choose a leader, how to handle concurrent proposals, how leader changes work. This led to every Paxos implementation being subtly different. Raft was explicitly designed to be understandable: it specifies every detail of the protocol, including the exact state machine of each node, the randomized election timeout mechanism, and the log matching property. In practice, Raft implementations (etcd, CockroachDB, TiKV, HashiCorp Consul) are more consistent with each other than Paxos implementations. Google still uses Paxos internally (Chubby lock service, Bigtable), but open-source ecosystem has largely standardized on Raft.
\end{curiobox}

\begin{curiobox}{What are CRDTs and when should you use them instead of Raft?}
CRDTs (Conflict-free Replicated Data Types) are data structures designed so that concurrent writes from different replicas can always be merged without conflicts. Counter CRDTs (G-Counter, PN-Counter) allow concurrent increments that are merged by summing per-replica counts. Set CRDTs allow elements to be added without coordination. The key property: for data types with natural merge functions, you never need coordination at write time, and convergence is guaranteed by the data structure itself. This is more efficient than Raft for data that fits the CRDT model, because Raft requires a round-trip to a leader for every write. The limitation: CRDTs only work for specific data types and access patterns. You cannot implement a financial ledger (where debits must not exceed balance) as a CRDT without giving up the no-overdraft guarantee. Riak (a distributed database) built its entire storage model around CRDTs; modern systems like Redis CRDT, Apple's CloudKit, and Amazon DynamoDB's convergence model use CRDT principles. Use CRDTs when: data is naturally additive, merge is associative and commutative, and eventual consistency is acceptable. Use Raft when: you need strict ordering, exclusion, or invariants that cannot be expressed as a CRDT.
\end{curiobox}

\begin{curiobox}{How does Jepsen testing find bugs in distributed databases?}
Jepsen is a testing framework created by Kyle Kingsbury (aphyr) that finds safety and liveness violations in distributed databases by injecting network partitions, clock skew, and process crashes while running concurrent workloads. A Jepsen test records all operations (their start time, end time, and result) and then checks the history for violations: did any read return a value that was never written? Did a linearizable system return a stale value after a write was acknowledged? Kingsbury's analyses (published since 2013) have found bugs in Cassandra, MongoDB, VoltDB, CockroachDB, etcd, and dozens of other systems --- bugs where the documentation claimed linearizability but the implementation silently violated it under specific failure conditions. The Jepsen reports are required reading for anyone designing or evaluating a distributed database. The practical lesson: never trust a database's consistency claims without looking for Jepsen results. As of 2024, Jepsen analyses have been published for over 30 databases and distributed systems.
\end{curiobox}

\begin{curiobox}{Why does network partition tolerance matter less in a single data center?}
In a well-run data center, network reliability is extremely high. Google's Borg cluster scheduler paper (2015) notes that they see approximately one network partition per year in a production cluster. However, individual machine failures are much more common: a cluster of 10,000 machines with a 0.1\% daily failure rate sees 10 machine failures per day. The distinction matters: a Raft cluster with 5 nodes can tolerate 2 simultaneous failures. If the 2 failures are in the form of 2 dead machines (which is common), Raft continues to operate normally. If the failure is a network partition that splits the 5 nodes into groups of 2 and 3 (which is rare inside a single DC), Raft still operates: the group of 3 has a majority and elects a leader. What Raft cannot handle is a partition that creates two groups of 2 and one group of 1 --- which requires at least 5 nodes going 3 different ways simultaneously, extremely unlikely in practice. The ``partition tolerance'' of CAP is most meaningful across data centers with less reliable WAN links.
\end{curiobox}

% ─────────────────────────────────────────────────────────────────────────────
\section{Self-Quiz}
% ─────────────────────────────────────────────────────────────────────────────

\begin{quizbox}
\textbf{Answer these as if the interviewer is asking. Cover your answers before attempting.}

\vspace{0.5em}

\textbf{Q1.} ``You've said you'll use Cassandra for this component. Walk me through what actually happens to the data during a network partition between your two data centers.''

\vspace{0.8em}

\textbf{Q2.} ``Your notification system needs to guarantee that every notification is delivered exactly once. How does your consistency model choice affect your ability to guarantee this?''

\vspace{0.8em}

\textbf{Q3.} ``You've proposed using a Raft cluster for leader election. Your cluster has 5 nodes. How many nodes can fail simultaneously before your system stops accepting writes, and why?''

\vspace{0.8em}

\textbf{Q4.} ``I've heard of vector clocks but never seen them in a real system. Can you give me a concrete example of what problem they solve that a simple timestamp doesn't?''

\vspace{0.8em}

\textbf{Q5.} ``Why don't you just use two-phase commit across your database shards to guarantee ACID transactions? That's the traditional approach.''
\end{quizbox}

\newpage

\begin{answerbox}
\textbf{Q1 Model Answer:}

With Cassandra's default configuration (quorum reads and writes with a replication factor of 3 across 2 DCs), a partition between the two data centers creates two isolated groups of nodes. Cassandra continues accepting writes on both sides. Each DC writes to its local nodes using the LOCAL\_QUORUM consistency level. When the partition heals, Cassandra applies \textit{read repair} (correcting stale replicas during reads) and \textit{hinted handoff} (replaying writes that the remote DC missed during the partition). The data converges to the most recent value, using timestamps to determine which write wins. The tradeoff: during the partition, a client reading from DC1 will not see writes that happened on DC2 during the partition, and vice versa. This is the availability-over-consistency tradeoff --- Cassandra stays up and accepting writes, but reads may be briefly stale. If that is unacceptable, you need a CP system.

\vspace{0.6em}

\textbf{Q2 Model Answer:}

Exactly-once delivery requires two things: deduplication on the consumer side (to handle at-least-once delivery retries) and an idempotency mechanism so that processing the same notification twice produces the same result. The consistency model affects the first part. If you use an eventually consistent store for deduplication state (``has notification ID X been sent?''), a race condition can occur: two consumer nodes check simultaneously, both see that X has not been sent, and both send it. To prevent this, the deduplication store must be linearizable --- you need a system (Redis with SETNX, DynamoDB with conditional writes, etcd) that guarantees exactly one winner when two processes attempt to set the same deduplication key simultaneously. The notification body can be stored in an AP system; only the deduplication state needs CP guarantees.

\vspace{0.6em}

\textbf{Q3 Model Answer:}

A 5-node Raft cluster can tolerate 2 simultaneous node failures. Raft requires a majority --- more than half the nodes --- to elect a leader and commit writes. A majority of 5 is 3. With 2 failures, 3 nodes remain, which is exactly the majority. If 3 nodes fail, only 2 remain, which is less than the majority of 5, and the cluster cannot elect a leader or commit writes. The formula is: a cluster of $N$ nodes tolerates $\lfloor (N-1)/2 \rfloor$ failures. For $N=5$, that is $\lfloor 4/2 \rfloor = 2$. This is why etcd recommends 3 or 5 nodes: 3 tolerates 1 failure, 5 tolerates 2 failures. 4 nodes would only tolerate 1 failure (same as 3 nodes) at higher cost, which is why even-numbered clusters are generally avoided.

\vspace{0.6em}

\textbf{Q4 Model Answer:}

Concrete example: Two users simultaneously add items to the same shopping cart from different regions. User adds Item A from US-East (timestamp 10:00:00.001) and Item B from EU-West (timestamp 10:00:00.001). The timestamps are identical due to clock drift. A timestamp-based system cannot determine which write should win or whether they are concurrent. With vector clocks: the US-East node's vector clock shows $[1, 0]$ (US-East incremented its counter), and the EU-West node's vector clock shows $[0, 1]$ (EU-West incremented its counter). Neither vector dominates the other, so the system detects that these two writes are concurrent --- they happened without knowledge of each other. The conflict resolution strategy for a shopping cart is to take the union: keep both items. If the vectors showed a causal relationship ($[2, 0]$ and $[1, 0]$, meaning the EU-West write knew about the US-East write), the later version would clearly supersede the earlier one. Vector clocks give you correct causality detection; timestamps give you unreliable causality inference.

\vspace{0.6em}

\textbf{Q5 Model Answer:}

Two-phase commit has a fundamental blocking problem: if the coordinator fails after all participants have completed Phase 1 (prepared) but before Phase 2 (commit or abort) begins, all participants are stuck holding locks indefinitely, waiting for the coordinator to recover. In a high-throughput distributed system, this creates cascading latency and timeout issues. Modern distributed databases solve this differently. CockroachDB uses Raft consensus within each shard, so each shard has a highly available leader without a single-point-of-failure coordinator. For cross-shard transactions, it uses a transaction record and MVCC timestamps rather than 2PC locking. Google Spanner does use 2PC across Paxos groups, but replaces each 2PC participant with a Paxos group so coordinator failure is handled by Paxos recovery without blocking. The short answer: 2PC is fine as a protocol, but it requires every participant and coordinator to be highly available. In practice, replacing single nodes with consensus groups at each 2PC role is expensive. Simpler and more scalable approaches like Raft-per-shard or Percolator-style MVCC are preferred for new systems.
\end{answerbox}

% ─────────────────────────────────────────────────────────────────────────────
\section{What's Next: L1-11 --- API Design}
% ─────────────────────────────────────────────────────────────────────────────

\begin{insightbox}
\textbf{Next lesson: L1-11 --- API Design (REST, gRPC, GraphQL, and the Interface Contracts That Hold Systems Together)}

You now have a mental model of how distributed nodes agree on data. The next question is: how do clients and servers agree on how to talk to each other? API design is the interface layer that sits on top of everything you have learned. In an interview, an interviewer who hears you correctly identify your consistency model and database choice will immediately follow up with ``Walk me through your API. What does the request look like?'' --- because the API is where your distributed system contracts become visible to the outside world.

L1-11 covers when to use REST versus gRPC versus GraphQL (each with real company examples of who uses what and why), how to version APIs without breaking clients, how streaming and webhooks differ from request-response APIs, and how to design idempotent endpoints --- which connects directly to the exactly-once delivery problem you encountered in this lesson's quiz question. After L1-11, you will have covered the full stack from networking (L1-01) through delivery semantics (L1-02) through distributed data guarantees (L1-10) through API contracts (L1-11) --- which together are sufficient to design almost any backend system an interviewer would ask about at the ICT3 level.

\textbf{Gate 1 progress:} L1-01 \checkmark{} + L1-02 \checkmark{} + L1-08 $\square$ → complete L1-08 (Message Queues) to unlock the Notification System retry mock interview.
\end{insightbox}

\end{document}
